% D.61; independent root full proof audit passed.
\subsection{Density of the abelian twists with zero Alexander cohomology}

\begin{statement}{D.61}{Proved}{Zero-rank Alexander modules occur in every completed abelian neighborhood}
For every pair \(a_0,b_0\in H^1(\mathbb Q,\mathbb Z_5(1))\)
and every integer \(N\ge0\), there is
\[
 a\in a_0+5^N H^1(\mathbb Q,\mathbb Z_5(1))
 \quad\hbox{such that}\quad
 H^1_{\mathrm{cts}}(G_{\mathbb Q},A_{(a,b_0)})=0.
                                                        \tag{D.61.1}
\]
The class \(a\) can be chosen with infinitely many nonzero
prime valuations, while the second coordinate is unchanged.
The abelian pairs with zero first Alexander cohomology are
therefore dense for the displayed five-adic neighborhood
system. Existence of a metabelian section over a general
pair is not asserted; if the fibre is nonempty, it is a singleton.
For \(b_0=0\), the actual cusp tangent supplies that singleton
in the geometrically pro-five metabelian quotient \(M\) of
D.35. No uniqueness assertion about the full all-prime
metabelian quotient is made.
\end{statement}

\begin{proof}
Choose actual continuous integral representatives of \(a_0,b_0\),
and let \(a_{0,q},b_{0,q}\in\mathbb Z_5\) denote their tame valuations
at \(q\). The family of valuations tends to zero five-adically:
for each \(h\), only finitely many primes have
either \(a_{0,q}\) or \(b_{0,q}\) outside \(5^h\mathbb Z_5\).
Choose \(q_0\ne5\) with
\(a_{0,q_0}\in5^{N+1}\mathbb Z_5\), and begin with
\(a_0+5^N\kappa(q_0)\). Its valuation at \(q_0\) is
\[
                c=a_{0,q_0}+5^N
                   \in5^N\mathbb Z_5^\times,
                                                        \tag{D.61.2}
\]
where the notation means \(v_5(c)=N\).
We never select \(q_0\) again.

At stage \(r\ge1\), put
\[
 \begin{split}
 a_{<r}&=a_0+5^N\kappa(q_0)+
             \sum_{j<r}5^{N+j}\sum_{q\in Q_j}\kappa(q),\\
 S_r&=\{5,\infty,q_0\}\cup\bigcup_{j<r}Q_j
       \cup\{q\ne5:\min(v_5(a_{0,q}),v_5(b_{0,q}))<r-1\}.
 \end{split}                                           \tag{D.61.3}
\]
The set \(S_r\) is finite, although \(a_0\) itself may
have unrestricted ramification. The action \(\rho_r\)
on \(W_r=\Lambda C/\mathfrak m^r\Lambda C\) determined
by \((a_{<r},b_0)\) is unramified outside \(S_r\), by
the binomial estimate (D.47.4). All additions at stages
\(j\ge r\) multiply it by
\(t_1^{5^{N+j}d}\equiv1\pmod{\mathfrak m^r}\), for
every integral \(d\). Thus its equality with the final
reduced action is exact, not merely up to conjugacy.

Enumerate the finite set
\(Z^1(G_{\mathbb Q,S_r},W_r)\). For each member
nonzero on \(\ker\rho_r\), choose its Chebotarev
witness prime exactly as in (D.60.6)--(D.60.7), with
\(\rho_r(\operatorname{Frob}_q)=1\) and nonzero
cocycle value. In addition to \(S_r\) and all earlier
chosen primes, require the witness to avoid the finite set
\[
 \{q:a_{0,q}\notin5^{N+r+1}\mathbb Z_5
                  \ \hbox{or}\ b_{0,q}\notin5^{N+r+1}\mathbb Z_5\}.
                                                        \tag{D.61.4}
\]
Removing finitely many primes does not affect existence
in a nonempty Chebotarev class. Let \(Q_r\) consist of
these distinct witness primes and one additional unused
prime satisfying the same avoidance condition. Define
\[
 a=a_0+5^N\kappa(q_0)+
                  \sum_{r\ge1}5^{N+r}\sum_{q\in Q_r}\kappa(q).
                                                        \tag{D.61.5}
\]
Uniform convergence is the same integral argument as
in D.60, and the entire difference is \(5^N\) times
an actual continuous integral cocycle. At every new
prime \(q\in Q_r\), the final valuation is
\[
               a_q=a_{0,q}+5^{N+r},
                    \qquad v_5(a_q)=N+r.
                                                        \tag{D.61.6}
\]
It cannot cancel: the first summand has strictly greater
five-adic valuation by (D.61.4). Later additions have
different prime supports. The additional prime at every
stage consequently ensures infinitely many nonzero
valuations, even if some original valuations of \(a_0\)
were already nonzero.

Every integral cocycle for the final \(A_{(a,b_0)}\),
reduced modulo \(\mathfrak m^r\), lies in the finite
set enumerated at stage \(r\). Indeed, at any prime
outside \(S_r\), both original valuations are divisible
by \(5^{r-1}\); if the prime is selected later, the
new summand is divisible by \(5^{N+j}\) with \(j\ge r\).
Both final valuations are therefore still divisible by
\(5^{r-1}\). D.43 makes a nonzero-inertia local
cocycle an integral coboundary, whose inertia value is
in \(\mathfrak m^r\) by (D.47.4). A zero-inertia
prime contributes zero by D.47. This proves the full
unrestricted cocycle capture, including primes already
ramified in \(a_0\).

The new nonzero inertia at each witness prime now
contradicts its prescribed nonzero Frobenius value in
exactly the manner of D.60. Hence every integral cocycle
descends continuously to the actual compact operator
image \(I\). The image of inertia at \(q_0\) is
\[
 N=\{\text{multiplication by }
         (t_1^c t_2^{b_{0,q_0}})^d:d\in\mathbb Z_5\},
                                                        \tag{D.61.7}
\]
where \(c\ne0\) is (D.61.2). Every actual operator is
a multiplier followed by \(t_i\mapsto t_i^{\chi_5(g)}\),
so it conjugates the displayed multiplier to its
\(\chi_5(g)\)-th power. Hence \(N\) is normal in \(I\).
Primitivity is unnecessary: \(t_1^c t_2^{b_{0,q_0}}-1\)
is nonzero, as is already seen by setting \(t_2=1\),
and is therefore a nonzero divisor in \(\Lambda\).
Thus \(A_{(a,b_0)}^N=0\). D.43 at \(q_0\), followed by the
normal-subgroup cocycle identity of D.60, proves
\(H^1(I,A_{(a,b_0)})=0\). This gives (D.61.1).

The exact D.35 torsor description makes every nonempty
metabelian fibre over the constructed pair a singleton.
For \(b_0=0\), the full cusp section with five-adic
tangent component \(a\) supplies that metabelian lift.
The topology in the statement is explicitly the one
with neighborhoods \(a_0+5^N H^1\); no compactness
of this completed Kummer group is asserted.
\end{proof}
